\begin{center}
    {\LARGE \textbf{2006分析化学}} \\[2ex]
    {\large 时量180分钟 \quad 满分90分}
\end{center}

\vspace{0.5cm}
\noindent\textbf{注意事项:}
\begin{enumerate}[label=\arabic*., parsep=0pt, itemsep=0pt]
    \item 答题前,考生先将自己的姓名、学号、考场号、座位号填写在答题卡上,将条形码准确粘贴在考生信息条形码粘贴区。
    \item 考试结束后,将本试卷与答题纸一并收回。
    \item 回答各个问题时,将答案写在答题纸上,写在本试卷上无效。
\end{enumerate}

\vspace{0.5cm}
\noindent\textbf{一、简答题（32 pts.）}

\begin{enumerate}[label=\arabic*., itemsep=1.5ex]
    \item 写出碘量法测定铜所涉及的三个反应式。
    \item 某同学如下配制$0.02\ \mathrm{mol/L}\ \ce{KMnO4}$溶液，请指出其错误。准确称取$3.161\ \mathrm{g}$固体$\ce{KMnO4}$，用煮沸过的去离子水溶解，转移至$1000\ \mathrm{ml}$容量瓶，稀释至刻度，然后用干燥的滤纸过滤。
    \item 下列情况对分析结果的影响（填偏高、偏低或无影响）
    \begin{enumerate}[label=(\arabic*), noitemsep, topsep=0pt, partopsep=0pt]
        \item $\ce{NaOH}$溶液浓度标定后由于保存不当吸收了$\ce{CO2}$，以标准溶液测定草酸摩尔质量。
        \item $\ce{NaOH}$溶液浓度标定后由于保存不当吸收了$\ce{CO2}$，若以此标准溶液测定$\ce{H3PO4}$浓度（甲基橙指示剂）。
    \end{enumerate}
    \item 用浓度为$0.1\ \mathrm{mol/L}$的氢氧化钠溶液滴定同浓度的柠檬酸溶液时，有几个滴定突跃？用什么指示剂？（已知柠檬酸的解离常数为$\mathrm{p}K_{\mathrm{a1}}=3.13$，$\mathrm{p}K_{\mathrm{a2}}=4.76$，$\mathrm{p}K_{\mathrm{a3}}=6.40$）
    \item 某采用以下方法测定溶液中$\ce{Ca^{2+}}$的浓度，请指出其错误。先用自来水洗净$250\ \mathrm{mL}$锥形瓶，用$50\mathrm{mL}$容量瓶量取试液倒入锥形瓶中，然后加少许铬黑T指示剂，用EDTA标准溶液滴定。
    \item 在采用EDTA配位滴定法测定金属离子时为什么要使用缓冲溶液？
    \item 影响沉淀溶解度的主要因素是什么？
    \item 写出下列溶液在化学计量点
    \begin{enumerate}[label=(\arabic*), noitemsep, topsep=0pt, partopsep=0pt]
        \item 用$\ce{NaOH}$滴定$\alpha-$氯乙酸（$\ce{HA}$，$\mathrm{p}K_\mathrm{a}=1.3$）和$\ce{NH4Cl}$混合溶液中的$\alpha-$氯乙酸至化学计量点时；
        \item 含$0.10\ \mathrm{mol/L}\ \ce{HCl}$和$0.20\ \mathrm{mol/L}\ \ce{H2SO4}$的混合溶液。
    \end{enumerate}
\end{enumerate}

\vspace{0.5cm}
\noindent\textbf{二、计算题（58 pts. total）}

\begin{enumerate}[label=\arabic*., itemsep=1.5ex]
    \item (10 pts.) 为了防止$1.0\times10^{-4}\ \mathrm{mol/L}\ \ce{Pb^{2+}}$在通$\ce{H2S}$时生成$\ce{PbS}$沉淀，溶液中$\ce{H+}$应控制为多少？（已知$K_{\mathrm{sp}}(\ce{PbS})=8\times10^{-28}$，$\ce{H2S}$的$K_{\mathrm{a1}}=5.7\times10^{-8}$，$K_{\mathrm{a2}}=1.2\times10^{-15}$，饱和$\ce{H2S}$浓度约为$0.10\ \mathrm{mol/L}$）
    \item (10 pts.) 已知二乙胺是二元胺盐（以$\ce{H2en^{2+}}$表示）的$\mathrm{p}K_{\mathrm{a1}}=6.85$，$\mathrm{p}K_{\mathrm{a2}}=9.93$，欲配制$1\ \mathrm{L}\ \mathrm{pH}=6.55$，总浓度为$0.1\ \mathrm{mol/L}$的二乙胺盐酸盐缓冲溶液，需加入固体$\ce{NaOH}$多少克？（已知$Mr(\ce{NaOH})=40.00$）
    \item (10 pts.) 在$1\ \mathrm{mol/L}\ \ce{HCl}$酸度下，$\ce{AsO4^{3-}}$能否氧化$\ce{I-}$析出$\ce{I2}$，而在$\mathrm{pH}=8$时$\ce{I2}$能否滴定$\ce{AsO3^{3-}}$生成$\ce{AsO4^{3-}}$，通过计算说明原因？（已知$E^\ominus_{\ce{AsO4/AsO3}}=0.56\ \mathrm{V}$，$E^\ominus_{\ce{I2/I-}}=0.54\ \mathrm{V}$）
    \item (10 pts.) 用碘量法测定某铜合金中铜的质量分数$w(\ce{Cu}\%)$，6次测定结果如下：$60.60$，$60.64$，$60.58$，$60.65$，$60.57$和$60.32$。(1) 用格鲁布斯法检测有无应合并的异常值（显著水平0.05）；(2) 估计铜的质量分数范围（$P=95\%$）；(3) 如果铜的质量分数标准值为$60.58\%$，试问测定有无系统误差（显著水平0.05）？
    \\
    已知：
    \begin{center}
    \begin{minipage}{0.45\linewidth}
    \centering
    \begin{tabular}{c|ccc}
    $n$ & 4 & 5 & 6 \\
    \hline
    $T_{0.05}$ & 1.463 & 1.672 & 1.832 \\
    \end{tabular}
    \end{minipage}
    \hfill
    \begin{minipage}{0.45\linewidth}
    \centering
    \begin{tabular}{c|ccc}
    $f$ & 4 & 5 & 6 \\
    \hline
    $t_{0.05}$ & 2.776 & 2.571 & 2.447 \\
    \end{tabular}
    \end{minipage}
    \end{center}
    \item (8 pts.) 试证明用沉淀掩蔽法在$\mathrm{pH}=12$时用EDTA能否准确滴定$\ce{Ca^{2+}}$、$\ce{Mg^{2+}}$混合溶液中的$\ce{Ca^{2+}}$而$\ce{Mg^{2+}}$不干扰。（已知$\ce{Ca^{2+}}$、$\ce{Mg^{2+}}$及EDTA浓度均为$0.020\ \mathrm{mol/L}$，$\ce{Mg(OH)2}$的$K_{\mathrm{sp}}=10^{-10.7}$，$\lg K_{\mathrm{CaY}}=10.7$，$\lg K_{\mathrm{MgY}}=8.7$，$\mathrm{pH}=12$时$\lg\alpha_{\mathrm{Y(H)}}=0$）
    \item (10 pts.) 碘在某有机溶剂和水之间的分配比$D$为$8.00$，如果$50.0\ \mathrm{ml}\ 0.100\ \mathrm{mol/L}\ \ce{I2}$水溶液用$100\ \mathrm{ml}$该有机溶剂振摇，问平衡后与$10.0\ \mathrm{ml}$有机相中$\ce{I2}$完全反应需要多少$0.0600\ \mathrm{mol/L}\ \ce{Na2S2O3}$溶液多少毫升？（已知$A_r(\ce{I})=127.0$）
\end{enumerate}